\documentclass[11pt]{article}

% Packages
\usepackage[margin=1in]{geometry}
\usepackage{amsmath}
\usepackage{amssymb}
\usepackage{graphicx}
\usepackage{booktabs}
\usepackage{siunitx}
\usepackage{hyperref}
\usepackage{xcolor}
\usepackage{float}

% Formatting
\setlength{\parindent}{0in}
\setlength{\parskip}{0.5\baselineskip}

\begin{document}

\begin{center}
  \Large\textbf{TA Solution \& Reference Guide} \\
  \vspace{0.1in}
  \huge\textbf{ENGR 212 Lab 7 - Active Filters}
\end{center}

\vspace{0.1in}
\hrule
\vspace{0.2in}

\section*{Lab Parameters \& Setup}
\begin{itemize}
  \item \textbf{Op-Amp:} LM324 Operational Amplifier.
  \item \textbf{Power Supply:} Split supply with $V_{CC} = \SI{+10}{\volt}$ and $V_{EE} = \SI{-10}{\volt}$.
  \item \textbf{Input Signal:} $\SI{1}{\volt}_{pp}$ sine wave from the function generator.
  \item \textbf{Components:} $R_1 = \SI{11}{\kilo\ohm}$, $R_2 = \SI{33}{\kilo\ohm}$, $C = \SI{0.01}{\micro\farad}$.
  \item \textbf{Measurement:} AC Voltage (RMS) using the Agilent 34401A DMM.
\end{itemize}

\vspace{0.2in}
% -------------------------------------------------------------------
\section*{Part 1: Active Low-Pass Filter}

\textbf{Theoretical Equations:}
\begin{itemize}
  \item \textbf{Transfer Function:} $H(s) = -\frac{R_2}{R_1} \frac{1}{1 + sR_2C}$
  \item \textbf{Cutoff Frequency ($f_c$):} $f_c = \frac{1}{2\pi R_2 C}$
  \item \textbf{Passband Gain (DC):} $|K| = \frac{R_2}{R_1}$
\end{itemize}

\textbf{Step-by-Step Calculations:}

\textbf{1. Cutoff Frequency ($f_c$):}
\begin{align*}
  f_c &= \frac{1}{2\pi R_2 C} \\
  f_c &= \frac{1}{2\pi (\SI{33e3}{\ohm})(\SI{0.01e-6}{\farad})} \\
  f_c &= \frac{1}{2\pi (\SI{3.3e-4}{\second})} \\
  f_c &= \frac{1}{0.002073} \\
  f_c &\approx \mathbf{482.3 \text{ \si{\hertz}}}
\end{align*}

\textbf{2. Passband Gain (Linear \& Decibels):}
\begin{align*}
  |K| &= \frac{R_2}{R_1} \\
  |K| &= \frac{\SI{33}{\kilo\ohm}}{\SI{11}{\kilo\ohm}} = \mathbf{3 \text{ V/V}} \\[1ex]
  |K|_{\text{dB}} &= 20 \log_{10}(|K|) \\
  |K|_{\text{dB}} &= 20 \log_{10}(3) \\
  |K|_{\text{dB}} &\approx 20 (0.4771) \approx \mathbf{9.54 \text{ dB}}
\end{align*}

\textbf{3. Gain at Cutoff Frequency:}
\begin{align*}
  \text{Gain at } f_c &= |K|_{\text{dB}} - 3\text{ dB} \\
  \text{Gain at } f_c &= 9.54\text{ dB} - 3\text{ dB} \\
  \text{Gain at } f_c &= \mathbf{6.54 \text{ dB}}
\end{align*}

\textit{TA Note on Data Collection:} Students are required to measure at $f_c$. Make sure they explicitly set their function generator to $\sim \SI{482}{\hertz}$ to take this measurement.

\vspace{0.2in}
% -------------------------------------------------------------------
\section*{Part 2: Active High-Pass Filter}

\textbf{Theoretical Equations:}
\begin{itemize}
  \item \textbf{Transfer Function:} $H(s) = -\frac{R_2}{R_1} \frac{sR_1C}{1 + sR_1C}$
  \item \textbf{Cutoff Frequency ($f_c$):} $f_c = \frac{1}{2\pi R_1 C}$
  \item \textbf{Passband Gain (High Frequency):} $|K| = \frac{R_2}{R_1}$
\end{itemize}

\textbf{Step-by-Step Calculations:}

\textbf{1. Cutoff Frequency ($f_c$):}
\begin{align*}
  f_c &= \frac{1}{2\pi R_1 C} \\
  f_c &= \frac{1}{2\pi (\SI{11e3}{\ohm})(\SI{0.01e-6}{\farad})} \\
  f_c &= \frac{1}{2\pi (\SI{1.1e-4}{\second})} \\
  f_c &= \frac{1}{0.000691} \\
  f_c &\approx \mathbf{1446.9 \text{ \si{\hertz}}}
\end{align*}

\textbf{2. Passband Gain (Linear \& Decibels):}
\begin{align*}
  |K| &= \frac{R_2}{R_1} \\
  |K| &= \frac{\SI{33}{\kilo\ohm}}{\SI{11}{\kilo\ohm}} = \mathbf{3 \text{ V/V}} \\[1ex]
  |K|_{\text{dB}} &= 20 \log_{10}(|K|) \\
  |K|_{\text{dB}} &= 20 \log_{10}(3) \approx \mathbf{9.54 \text{ dB}}
\end{align*}

\textbf{3. Gain at Cutoff Frequency:}
\begin{align*}
  \text{Gain at } f_c &= |K|_{\text{dB}} - 3\text{ dB} \\
  \text{Gain at } f_c &= 9.54\text{ dB} - 3\text{ dB} \\
  \text{Gain at } f_c &= \mathbf{6.54 \text{ dB}}
\end{align*}

\textit{TA Note on Data Collection:} Make sure students sweep up to at least $\SI{10000}{\hertz}$ to clearly see the plateau of the passband at $\sim 9.54 \text{ dB}$.

\newpage
% -------------------------------------------------------------------
\section*{Common Bench Troubleshooting \& Pitfalls}

If a student's circuit is outputting \SI{0}{\volt}, clipping heavily, or yielding completely flat frequency responses, check the following:

\begin{enumerate}
  \item \textbf{Function Generator ``High-Z'' Mismatch:} By default, Agilent function generators assume a $\SI{50}{\ohm}$ load. If a student sets it to $\SI{1}{\volt}_{pp}$, the open-circuit voltage will actually be $\SI{2}{\volt}_{pp}$. \textbf{Fix:} The lab instructions correctly tell students to ``Measure both $V_{IN}$ and $V_{OUT}$ with the DMM''. As long as they calculate gain using their \textit{actual} DMM $V_{IN}$ measurement rather than assuming the function generator display is correct, their dB calculations will be perfect.

  \item \textbf{DMM AC vs. DC Mode:} Students must use the \textbf{AC Voltage} setting on the DMM. If they measure in DC, they will read $\sim \SI{0}{\volt}$ since the average of a sine wave is zero.

  \item \textbf{LM324 Power Rails:}
    \begin{itemize}
      \item $V_{CC}$ ($\SI{+10}{\volt}$) must be connected to \textbf{Pin 4}.
      \item $V_{EE}$ ($\SI{-10}{\volt}$) must be connected to \textbf{Pin 11}.
      \item \textit{Frequent Mistake:} Students tie Pin 11 to Ground instead of the negative power supply, which will severely clip the negative half of the AC sine wave.
    \end{itemize}

  \item \textbf{Breadboard Parasitics at \SI{10}{\kilo\hertz}:} At higher frequencies, loose wires and breadboard capacitance can slightly degrade the signal. If their $\SI{10}{\kilo\hertz}$ measurement droops slightly below $9.54 \text{ dB}$, this is a normal real-world non-ideality.
\end{enumerate}

\vspace{0.2in}
% -------------------------------------------------------------------
\section*{Grading the Formal Report}

Unlike Lab 6, this experiment requires a \textbf{Formal Lab Report}. Ensure students follow the standard formal structure (Abstract/Introduction, Methodology, Results, Discussion, Conclusion).

\textbf{What to look for in the Data \& Results:}
\begin{itemize}
  \item \textbf{Data Tables:} They must have calculated the theoretical gain alongside their measured $V_{OUT}/V_{IN}$. Gains must be properly converted to Decibels ($20 \log_{10}(V_{out}/V_{in})$).

  \item \textbf{Bode Plots (Magnitude Only):}
    \begin{itemize}
      \item The x-axis \textbf{must} be on a logarithmic scale (Semilog graph).
      \item Both theoretical and measured curves must be plotted on the \textit{same graph} for comparison.
      \item Low-Pass should start flat at $\sim 9.5 \text{ dB}$ and roll off at $-20 \text{ dB/decade}$ after $\SI{482}{\hertz}$.
      \item High-Pass should rise at $+20 \text{ dB/decade}$ and flatten out at $\sim 9.5 \text{ dB}$ after $\SI{1447}{\hertz}$.
    \end{itemize}

  \item \textbf{Error Analysis:} The manual specifically asks to quantify the percentage error between measured and calculated data \textit{at the cutoff frequencies}. Ensure they explicitly calculated error at $\SI{482}{\hertz}$ and $\SI{1447}{\hertz}$. Standard component tolerances (resistors $\pm 5\%$, capacitors $\pm 10\% - 20\%$) usually result in $f_c$ errors of around $5\% - 15\%$.
\end{itemize}

\end{document}
